\documentclass[11pt]{article}

\usepackage[a4paper,margin=28mm]{geometry}
\usepackage[T1]{fontenc}
\usepackage{lmodern}
\usepackage{microtype}
\usepackage{amsmath,amssymb,mathtools}
\usepackage{booktabs}
\usepackage{xcolor}
\usepackage[hidelinks]{hyperref}

\newcommand{\Q}{\mathbb{Q}}
\newcommand{\Z}{\mathbb{Z}}

\title{An Exact Rational Solution with \(y=9162\)}
\author{}
\date{12 July 2026}

\begin{document}
\maketitle

\begin{abstract}
This note records and verifies an exact rational nonintegral solution of a
Diophantine equation with two rational parameters.  The solution gives
\(y=9162\) (and, by symmetry, \(y=-9162\)).  All calculations are performed
in exact rational arithmetic.  A final remark explains why this solution is
valid for the original rational search but does not meet the later, stronger
restriction that the reduced denominators of \(k_1\) and \(k_2\) divide the
two large coefficients.
\end{abstract}

\section{The equation}

Set
\begin{align*}
A&=7729484335457653901640057298531371241781,\\
B&=7668575607239450973459863267707132263860,\\
C&=2486598372481845396683104279916570951657,\\
D&= 609530524018264138310326718615033307496.
\end{align*}
We consider
\begin{equation}\label{eq:original}
 y^2=
 \bigl(6(Ak_1+B)+(Ck_2+D)\bigr)^2
 +\frac{36(Ak_1+B)^3-19}{Ck_2+D}.
\end{equation}
It is convenient to introduce
\begin{equation}\label{eq:nxdef}
 n=Ak_1+B,\qquad x=Ck_2+D.
\end{equation}
Then \eqref{eq:original} becomes
\begin{equation}\label{eq:reduced}
 y^2=(6n+x)^2+\frac{36n^3-19}{x}.
\end{equation}

\section{The rational parameter values}

Take
\begin{equation}\label{eq:nxvalues}
 n=-\frac{1506}{5},\qquad x=-\frac{61}{5}.
\end{equation}
Solving \eqref{eq:nxdef} for \(k_1,k_2\) gives
\begin{equation}\label{eq:kformulas}
 k_1=\frac{n-B}{A},\qquad k_2=\frac{x-D}{C}.
\end{equation}
Substituting \eqref{eq:nxvalues} into \eqref{eq:kformulas} yields
\begin{equation}\label{eq:k1}
\boxed{
 k_1=-\frac{
 38342878036197254867299316338535661320806
 }{
 38647421677288269508200286492656856208905
 }}
\end{equation}
and
\begin{equation}\label{eq:k2}
\boxed{
 k_2=-\frac{
 3047652620091320691551633593075166537541
 }{
 12432991862409226983415521399582854758285
 }}.
\end{equation}
Indeed, the numerators and denominators arise from
\begin{align*}
5B+1506
 &=38342878036197254867299316338535661320806,\\
5A
 &=38647421677288269508200286492656856208905,\\
5D+61
 &=3047652620091320691551633593075166537541,\\
5C
 &=12432991862409226983415521399582854758285.
\end{align*}
Moreover,
\begin{align*}
\gcd(5B+1506,5A)&=1,\\
\gcd(5D+61,5C)&=1.
\end{align*}
Thus the fractions in \eqref{eq:k1} and \eqref{eq:k2} are reduced, and both
\(k_1\) and \(k_2\) are rational nonintegers.

\section{Exact verification of the affine substitutions}

Using \eqref{eq:k1},
\begin{align*}
Ak_1+B
&=A\left(-\frac{5B+1506}{5A}\right)+B\\
&=-\frac{5B+1506}{5}+B\\
&=-\frac{1506}{5}.
\end{align*}
Similarly, using \eqref{eq:k2},
\begin{align*}
Ck_2+D
&=C\left(-\frac{5D+61}{5C}\right)+D\\
&=-\frac{5D+61}{5}+D\\
&=-\frac{61}{5}.
\end{align*}
Consequently, the large-parameter equation \eqref{eq:original} is reduced
exactly to \eqref{eq:reduced} with the values in \eqref{eq:nxvalues}.

\section{Exact verification that \texorpdfstring{\(y=9162\)}{y=9162}}

First,
\begin{equation}\label{eq:linearpart}
6n+x
=6\left(-\frac{1506}{5}\right)-\frac{61}{5}
=-\frac{9097}{5},
\end{equation}
so
\begin{equation}\label{eq:squarepart}
(6n+x)^2=\frac{9097^2}{25}
=\frac{82755409}{25}.
\end{equation}
Next,
\begin{align*}
36n^3-19
&=36\left(-\frac{1506}{5}\right)^3-19\\
&=-\frac{122963842151}{125}.
\end{align*}
Since \(x=-61/5\), it follows that
\begin{align}
\frac{36n^3-19}{x}
&=\left(-\frac{122963842151}{125}\right)
  \left(-\frac{5}{61}\right)\notag\\
&=\frac{2015800691}{25}.\label{eq:cubicpart}
\end{align}
Combining \eqref{eq:squarepart} and \eqref{eq:cubicpart},
\begin{align*}
(6n+x)^2+\frac{36n^3-19}{x}
&=\frac{82755409+2015800691}{25}\\
&=\frac{2098556100}{25}\\
&=83942244\\
&=9162^2.
\end{align*}
Therefore the exact solutions for \(y\) are
\begin{equation}
\boxed{y=9162\quad\text{or}\quad y=-9162.}
\end{equation}
In particular, the requested positive value \(y=9162\) is valid.

\section{Status under the later denominator restriction}

The reduced denominators of the two rational parameters are
\begin{align*}
\operatorname{den}(k_1)&=5A
 =38647421677288269508200286492656856208905,\\
\operatorname{den}(k_2)&=5C
 =12432991862409226983415521399582854758285.
\end{align*}
Hence this solution is a valid solution of the original rational-parameter
problem, but it does \emph{not} satisfy the later conditions
\[
\operatorname{den}(k_1)\mid A,
\qquad
\operatorname{den}(k_2)\mid C.
\]
The obstruction is exactly the extra factor \(5\), which is also the common
denominator of \(n=-1506/5\) and \(x=-61/5\).

\clearpage
\appendix
\section{Exact Magma verification}

The following Magma code verifies the solution without floating-point
arithmetic.

{\small
\begin{verbatim}
Q := Rationals();

A := 7729484335457653901640057298531371241781;
B := 7668575607239450973459863267707132263860;
C := 2486598372481845396683104279916570951657;
D := 609530524018264138310326718615033307496;

k1 := Q!(-38342878036197254867299316338535661320806)
      / 38647421677288269508200286492656856208905;

k2 := Q!(-3047652620091320691551633593075166537541)
      / 12432991862409226983415521399582854758285;

y := 9162;

n := A*k1 + B;
x := C*k2 + D;

assert n eq Q!(-1506)/5;
assert x eq Q!(-61)/5;
assert y^2 eq (6*n + x)^2 + (36*n^3 - 19)/x;

assert GCD(Abs(Numerator(k1)), Denominator(k1)) eq 1;
assert GCD(Abs(Numerator(k2)), Denominator(k2)) eq 1;

print "n =", n;
print "x =", x;
print "y =", y;
print "RHS =", (6*n + x)^2 + (36*n^3 - 19)/x;
\end{verbatim}
}

\end{document}
